
\subsection{扰动后的模空间逸出仓西邻域 I}
\label{subsec:7-2-3}

本小节先指出构造相容多值截面族时
出现的另一困难；
它与注~3.4.10 提到的问题有关。
然后说明克服困难的思路，
详细做法留待后面几小节。

从经典分量
$\mathcal M^{\mathrm{main}}_{k+1}(\beta_0)$
的最简单情形开始；
处理其他 $\beta$ 的
$\mathcal M^{\mathrm{main}}_{k+1}(\beta)$
时也会遇到同一问题。
回忆
$\mathcal M^{\mathrm{main}}_3(\beta_0)$
的扰动可以识别为
\[
\left\{
\bigl(x,\varphi_0^{-1}(x),
\varphi_1^{-1}(x)\bigr)
\ \middle|\ x\in L
\right\},
\]
其中
$\varphi_0,\varphi_1\colon L\to L$
是充分接近恒等映射的微分同胚。
于是
$\mathcal M^{\mathrm{main}}_4(\beta_0)$
的扰动由一族路径
$\varphi_{3,t}\colon L\to L^3$
（$t\in[0,1]$）的并描述，
其两端满足
\[
\begin{aligned}
\varphi_{3,0}(x)
={}&\bigl(
\varphi_0^{-1}(\varphi_0^{-1}(x)),
\varphi_0^{-1}(\varphi_1^{-1}(x)),
\varphi_1^{-1}(\varphi_1^{-1}(x))
\bigr),\\
\varphi_{3,1}(x)
={}&\bigl(
\varphi_0^{-1}(\varphi_0^{-1}(x)),
\varphi_1^{-1}(\varphi_0^{-1}(x)),
\varphi_1^{-1}(\varphi_1^{-1}(x))
\bigr).
\end{aligned}
\]
也就是说，定义
\begin{equation}
\mathcal M^{\mathrm{main}}_4(\beta_0)^s
=\{(x,\varphi_{3,t}(x))\mid
x\in L,\ t\in[0,1]\}
\subseteq L^4.
\tag{7.2.24}\label{eq:7-2-24}
\end{equation}
这个集合仍在对角线的邻域内。
但是，
如果以 $\epsilon$ 估计
$\operatorname{dist}
(\varphi_0^{-1}(x),\varphi_1^{-1}(x))$，
那么
$\mathcal M^{\mathrm{main}}_4(\beta_0)^s$
到对角线的距离
只能用大约 $2\epsilon$ 估计。
随着 $k$ 增大，
$\mathcal M^{\mathrm{main}}_{k+1}(\beta_0)^s$
的估计会更差。
例如
$\mathcal M^{\mathrm{main}}_{k+2}(\beta_0)^s$
含有形如
\begin{equation}
\bigl(
x,\varphi_0^{-k}(x),
\varphi_1^{-1}(\varphi_0^{1-k}(x)),
\ldots,\varphi_1^{-k}(x)
\bigr)
\tag{7.2.25}\label{eq:7-2-25}
\end{equation}
的点，
它到对角线的距离约为 $k\epsilon$。
因此，$k\to\infty$ 时，
$\mathcal M^{\mathrm{main}}_{k+2}(\beta_0)^s$
可能逸出对角线的这个邻域，
也就是逸出为
$\mathcal M^{\mathrm{main}}_{k+2}(\beta_0)$
选定的仓西邻域。

在
$\mathcal M^{\mathrm{main}}_3(\beta_0)$
上选取扰动 $s$ 时，
可以把 $\epsilon$ 取得任意小，
从而对固定的 $k_0$，
半径 $k_0\epsilon$ 的邻域
仍包含在仓西邻域中。
可是必须先固定
$\mathcal M^{\mathrm{main}}_3(\beta_0)$
的扰动，
而在构造多值截面系统的过程中
不能改变它。
故即使 $\epsilon$ 取得很小，
$k\to\infty$ 时
$k\epsilon$ 仍可变得很大。

同时扰动无限多个模空间，
除上述困难外还有其他问题。
另一方面，
上面的论证也清楚表明：
只要适当选取 $s$，
对任意大但\emph{固定}的 $k$，
构造都可以继续。
换言之，
$\epsilon$ 要依赖于 $k$ 选取。
具体到
$\mathcal M^{\mathrm{main}}_{k+1}(\beta_0)$，
对任意给定的 $k$
都可以构造 $A_k$ 结构；
但要由这些结构构造
$A_\infty$ 结构，
还需要额外论证。
这一论证见第~7.2.8 小节的命题 7.2.93，
其证明将在第~7.2.11 小节给出。
类似问题曾在哈密顿不动点的
弗洛尔上同调中
由 \textbf{[Ono95]} 以另一方式处理。

第~7.2.6--7.2.8 小节
将以如下同调代数论证克服这一困难。
考虑定理~\ref{thm:3-4-8}
证明中出现的 $A_\infty$ 结构：
它由
$\mathcal M^{\mathrm{main}}_{k+1}(\beta_0)$
构造，不含量子效应。
第~7.2.2 小节概述、
第~7.2.4 和 7.2.5 小节详述的论证表明，
对每个固定的 $K$，
可以为全部 $k\le K$ 的
\[
\mathcal M^{\mathrm{main}}_{k+1}
(\beta_0;P_1,\ldots,P_k)
\]
选取扰动
（或障碍丛的截面），
使之满足适当相容条件。
第~7.2.7 小节将用这些扰动
定义一个 $A_K$ 结构，
记作 $\mathfrak m^{(K)}$。
由第~4.6 节的论证可知，
若 $K<K'$，
则将 $A_{K'}$ 结构
$\mathfrak m^{(K')}$
视为 $A_K$ 结构时，
它与 $A_K$ 结构
$\mathfrak m^{(K)}$
同伦等价。

随后将证明：
若 $A_K$ 代数
$\mathfrak m_k$（$k=1,\ldots,K$）
作为 $A_K$ 代数，
与某个 $A_{K+1}$ 代数
$\mathfrak m'_k$
同伦等价，
则存在 $\mathfrak m_{K+1}$，
使原有的
$\mathfrak m_1,\ldots,\mathfrak m_K$
连同 $\mathfrak m_{K+1}$
成为 $A_{K+1}$ 结构。
这就是后文定理 7.2.72；
其证明类似定理~\ref{thm:4-2-45}
的证明，并以障碍理论为基础。
这样便可得到 $A_\infty$ 结构。

实际上还要构造包含量子效应的
滤过 $A_\infty$ 代数。
由于 $\mathfrak m_0$ 的出现
使定义更加困难，
第~4.4.2 小节并未定义
滤过 $A_K$ 结构。
第~7.2.6 小节将定义
滤过 $A_{n,K}$ 结构，
用于滤过 $A_\infty$ 代数的归纳构造。
